\documentclass[
	12pt,				% tamanho da fonte
	openright,			% capítulos começam em pág ímpar (insere página vazia caso preciso)
	oneside,			% para impressão em verso e anverso. Oposto a oneside
	a4paper,			% tamanho do papel. 
	english,			% idioma adicional para hifenização
	french,				% idioma adicional para hifenização
	spanish,			% idioma adicional para hifenização
	brazil				% o último idioma é o principal do documento
	]{abntex2}
	
\usepackage{lmodern}				% Usa a fonte Latin Modern			
\usepackage[T1]{fontenc}			% Selecao de codigos de fonte.
\usepackage[utf8]{inputenc}			% Codificacao do documento (conversão automática dos acentos)
\usepackage{lastpage}				% Usado pela Ficha catalográfica
\usepackage{indentfirst}			% Indenta o primeiro parágrafo de cada seção.
\usepackage{color}					% Controle das cores
\usepackage{graphicx}				% Inclusão de gráficos e imagens
\usepackage{microtype}				% para melhorias de justificação
\usepackage{mathtools}				% Matemática
\usepackage{multirow}				% Tabelas
\usepackage[table,xcdraw]{xcolor}	% Tabelas (cores)
\usepackage{pdfpages}
\usepackage{lscape}
\usepackage{float}
\usepackage{imakeidx}
\usepackage{amsmath}
\usepackage{amsthm}
\usepackage{amssymb}
\usepackage{appendix}
\usepackage{enumerate}
\usepackage[pngtex]{pict2e}

\newtheorem{teorema}{Teorema}[chapter]
\newtheorem{corolário}[teorema]{Corolário}
\newtheorem{definição}[teorema]{Definição}
\newtheorem{exemplo}[teorema]{Exemplo}
\newtheorem{proposição}[teorema]{Proposição}
\newtheorem{observação}[teorema]{Observação}
\newtheorem{notação}[teorema]{Notação}
\newtheorem{lema}[teorema]{Lema}


\usepackage{amsfonts}
\usepackage{algorithm}
\usepackage[noend]{algpseudocode}
\usepackage{array}
%\usepackage{gensymb}
\usepackage{float}

\hyphenpenalty=10000

\usepackage{lipsum}				% para geração de dummy text

\usepackage[
    backend=biber,
    style=numeric,
  ]{biblatex}

 \usepackage{csquotes}

 \addbibresource{bibliography.bib}
%--------------------------------

\newcommand{\STAB}[1]{\begin{tabular}{@{}c@{}}#1\end{tabular}}
	
\titulo{Título da dissertação ou tese}
\autor{Nome do autor}
\local{São Luís - MA}
\data{2023}
\orientador{Prof. Dr. }
%\coorientador{Prof. Dr. Anselmo Cardoso de Paiva} %se houver

\tipotrabalho{Dissertação}
\preambulo{Dissertação apresentada como requisito parcial para obtenção do título de Mestre em Matemática, ao Programa de Pós-Graduação em Matemática, da Universidade Federal do Maranhão.}

% alterando o aspecto da cor azul
\definecolor{blue}{RGB}{41,5,195}

% informações do PDF
\makeatletter
\hypersetup{
     	%pagebackref=true,
		pdftitle={\@title}, 
		pdfauthor={\@author},
    	pdfsubject={\imprimirpreambulo},
	    pdfcreator={LaTeX with abnTeX2},
		pdfkeywords={abnt}{latex}{abntex}{abntex2}{trabalho acadêmico}, 
		colorlinks=true,       		% false: boxed links; true: colored links
    	linkcolor=blue,          	% color of internal links
    	citecolor=blue,        		% color of links to bibliography
    	filecolor=magenta,      		% color of file links
		urlcolor=blue,
		bookmarksdepth=4
}
\makeatother

% O tamanho do parágrafo é dado por:
\setlength{\parindent}{1.3cm}

% Controle do espaçamento entre um parágrafo e outro:
\setlength{\parskip}{0.2cm}  % tente também \onelineskip

\makeindex[columns=2, title=Índice Remissivo, intoc, 
           options= -s examplestyle.ist]

\begin{document}

\selectlanguage{brazil}

\frenchspacing 

% ---
% Capa
% ---
\renewcommand{\imprimircapa}{%
\begin{capa}%
\center
\includegraphics[scale=0.25]{imagens/logoufma.png} \\ 
\vspace{1cm}
\ABNTEXchapterfont\Large UNIVERSIDADE FEDERAL DO MARANHÃO \\
\ABNTEXsectionfont\Large Programa de Pós-Graduação em Matemática \\

\vspace*{5cm}

\begin{center}
{\ABNTEXsectionfont\Large\imprimirautor} \\
%\vfill
\ABNTEXchapterfont\bfseries\LARGE\imprimirtitulo
\end{center}
\vfill
\large\imprimirlocal \\
\large\imprimirdata
\vspace*{1cm}
\end{capa}
}	
\imprimircapa

%CASO PRECISE
%\includepdf[pages={1}]{capa_dissertacao.pdf}

\imprimirfolhaderosto


\includepdf[pages={1}]{ficha_esboço.pdf}
%\includepdf[pages={1}]{assinaturas.pdf}

\begin{folhadeaprovacao}

  \begin{center}
    {\ABNTEXchapterfont\large\imprimirautor}

    \vspace*{\fill}\vspace*{\fill}
    \begin{center}
      \ABNTEXchapterfont\bfseries\Large\imprimirtitulo
    \end{center}
    \vspace*{\fill}
    
    \hspace{.45\textwidth}
    \begin{minipage}{.5\textwidth}
        \imprimirpreambulo
    \end{minipage}%
    \vspace*{\fill}
  \end{center}
        
  Dissertação de Mestrado. \imprimirlocal, 12 de Maio de 
  2023:

  \assinatura{\textbf{\imprimirorientador} \\ Orientador\\ Universidade Federal do Maranhão} 
  %\assinatura{\textbf{\imprimircoorientador} \\ Coorientador\\ Universidade Federal do Maranhão} 
   
  \assinatura{\textbf{Prof. Dr. } \\ Examinador Interno \\ Universidade Federal do Maranhão}
  \assinatura{\textbf{Prof. Dr. } \\ Examinador Externo \\ Universidade de São Paulo - ICMC}
      
  \begin{center}
    \vspace*{0.5cm}
    {\large\imprimirlocal}
    \par
    {\large\imprimirdata}
    \vspace*{1cm}
  \end{center}
  
\end{folhadeaprovacao}

% ---
% Dedicatória
% ---
\begin{dedicatoria}
  \vspace*{\fill}
  %\RaggedRight
  %\noindent
  \begin{flushright}
        \textit{À pequenininha Alice.}
  \end{flushright}
  \vspace*{\fill}
\end{dedicatoria}%
% ---

% ---
% Agradecimentos
% ---
\begin{agradecimentos}
Agradeço 

\end{agradecimentos}
% ---

% ---
% Epígrafe
% ---
\newpage
\begin{epigrafe}
   \vspace*{\fill}
	\begin{flushright}
		\textit{"A matemática, vista corretamente, possui não 
  apenas verdade, mas também suprema beleza - uma beleza fria 
  e austera, como a da escultura."}

    (Bertrand Russell)
	\end{flushright}
\end{epigrafe}
% ---


% resumo em português
\setlength{\absparsep}{18pt} % ajusta o espaçamento dos parágrafos do resumo
\begin{resumo}

\end{resumo}

% resumo em inglês
\begin{resumo}[Abstract]
\begin{otherlanguage*}{english}

\end{otherlanguage*}
\end{resumo}

\cleardoublepage

\pdfbookmark[0]{\contentsname}{toc}
\tableofcontents*
\cleardoublepage



\input{1_introducao.tex}
\input{2_relacionados.tex}
\input{3_referencial.tex}
\input{appendix.tex}

\appendix

\phantompart

\postextual

\printbibliography

\printindex

\end{document}